% Figure 1 --- the two axes of the paper.
%   (a) how a monitor observes, and whether that observation enters the trajectory
%   (b) what happens after it fires, and how much true state survives
% Explicit coordinates are used throughout so the layout is predictable.
% Requires only tikz libraries already loaded in the preamble
% (arrows.meta, backgrounds, fit, positioning).
\begin{tikzpicture}[
    font=\small,
    turn/.style ={draw=agentblue!55, fill=agentblue!8, rounded corners=2pt,
                  minimum width=11mm, minimum height=6.5mm, inner sep=1pt, align=center,
                  font=\footnotesize},
    mon/.style  ={draw, rounded corners=2pt, minimum width=30mm, minimum height=10mm,
                  inner sep=2pt, align=center, font=\scriptsize},
    op/.style   ={draw, rounded corners=2pt, minimum width=44mm, minimum height=11mm,
                  inner sep=2pt, align=center, font=\scriptsize},
    panel/.style={font=\footnotesize\bfseries, anchor=west},
    note/.style ={font=\scriptsize\itshape, text=black!55, align=left},
    read/.style ={-{Stealth[length=1.8mm]}, thick, draw=black!45},
    flow/.style ={-{Stealth[length=1.6mm]}, draw=agentblue!60},
    write/.style={-{Stealth[length=1.8mm]}, thick, dashed, draw=sentinelorange}
  ]

  %% ============================== panel (a) : observation
  \node[panel] at (0,3.05) {(a)\, How the monitor observes};

  \foreach \i/\lbl in {0/$t$, 1/$t{+}1$, 2/$t{+}2$, 3/$t{+}3$}
    \node[turn] (t\i) at ({0.55+\i*1.35}, 2.30) {\lbl};
  \foreach \a/\b in {0/1, 1/2, 2/3} \draw[flow] (t\a) -- (t\b);
  \node[note] at (6.05,2.30) [anchor=west] {agent trajectory};

  \node[mon, draw=probeviolet, fill=probeviolet!8] (internal) at (2.6,3.55)
       {internal probe\\reads hidden activations};
  \draw[read, draw=probeviolet] (t1.north) -- (t1.north |- internal.south);

  \node[mon, draw=passiveteal, fill=passiveteal!8] (passive) at (1.35,0.95)
       {\textbf{passive} monitor\\reads the emitted trace};
  \node[mon, draw=sentinelorange, fill=sentinelorange!10] (active) at (5.05,0.95)
       {\textbf{active} monitor\\adds a chore to the context};

  \draw[read, draw=passiveteal]    (t1.south) -- (passive.north);
  \draw[read, draw=sentinelorange] (t2.south) -- (active.north);
  \draw[write] (active.east) -- (6.85,0.95) -- (6.85,2.30) -- (t3.east);

  \node[note, text=sentinelorange, text width=52mm] at (0.1,0.05) [anchor=north west]
       {only the active monitor writes back: the observation becomes part of the
        trajectory it is measuring};

  %% ============================== panel (b) : recovery
  \node[panel] at (0,-1.25) {(b)\, What happens when it fires};

  \node[mon, draw=warningred, fill=warningred!8, minimum width=24mm]
       (alarm) at (1.15,-2.25) {signal fires\\degradation suspected};

  \node[op, draw=warningred, fill=warningred!6]     (compact)  at (4.60,-1.65)
       {\textbf{compaction} (lossy)\\the \emph{agent} summarises its own state;\\
        errors it has already made survive};
  \node[op, draw=recoverygreen, fill=recoverygreen!8] (reground) at (4.60,-2.95)
       {\textbf{re-grounding} (deterministic)\\the \emph{harness} rebuilds state\\
        from an external record};

  \draw[read, draw=warningred]     (alarm.east) -- (2.75,-2.25) -- (2.75,-1.65) -- (compact.west);
  \draw[read, draw=recoverygreen]  (alarm.east) -- (2.75,-2.25) -- (2.75,-2.95) -- (reground.west);

  \node[turn, minimum width=14mm] (r1) at (8.05,-1.65) {resume};
  \node[turn, minimum width=14mm] (r2) at (8.05,-2.95) {resume};
  \draw[read, draw=warningred]    (compact.east)  -- (r1.west);
  \draw[read, draw=recoverygreen] (reground.east) -- (r2.west);

  \node[note, text width=90mm] at (0.1,-3.75) [anchor=north west]
       {The same monitor can help or hurt depending on which branch it triggers.};

  \begin{scope}[on background layer]
    \node[fit=(internal)(active)(t3), draw=black!15, rounded corners=3pt, inner sep=3.5mm] {};
    \node[fit=(alarm)(compact)(reground)(r1)(r2), draw=black!15, rounded corners=3pt,
          inner sep=3.5mm] {};
  \end{scope}
\end{tikzpicture}
